\section{Calendar-time optimality and dispatch verification}
\label{app:dispatch}

\subsection{Existence and calendar-horizon optimality}

We first record regularity and growth properties needed in both verification
arguments. The clipped loads are continuous and locally Lipschitz for
$q>0$. Thus $C$ and $J$ are continuous and locally Lipschitz there, including
across load and highway-capacity switches. Since $n_i\le q$, the stem terms
in the batch service cost imply
\begin{equation}
 C(q)\ge4c_H\rho>0,
 \qquad
 J(q)\ge\frac{4c_H\rho\Lambda}q+\frac{\omega q}2.
 \label{eq:proof-growth-lower}
\end{equation}
In particular $J$ diverges at zero and at infinity. It attains its minimum
$g$ on a compact subset of $(0,\infty)$, and its minimizer set
$\mathcal B$ is nonempty and compact. This argument does not require
convexity of the reduced objective or uniqueness of its minimizer.

There is also a constant $K_0<\infty$ such that
\begin{equation}
 0\le C(q)\le K_0(1+q),\qquad q>0.
 \label{eq:proof-growth-upper}
\end{equation}
For $q\ge1$, evaluate the unoptimized service cost at the fixed feasible
loads $n_1=n_2=\min(M,1)$. Then $s(q)\le a+b$, and the per-shipment
service cost is bounded. For $0<q\le1$, the feasible choice
$n_1=n_2=\min(M,q)$ gives a bounded batch cost after multiplication by
$q$: the stems are bounded, $qs(q)=aq+b\sqrt q$ is bounded, and
$q/\min(q,M_A)$ is bounded. Optimizing the loads can only reduce this
upper bound. In fact, the clipped optimizers equal $q$ for sufficiently
small $q$, and
\begin{equation}
 \lim_{q\downarrow0}C(q)=4c_H\rho+c_A\tau_{\mathrm{phys}}.
 \label{eq:proof-zero-limit}
\end{equation}

For a policy in Section~\ref{subsec:dynamic-wave}, let $r$ be the elapsed
time since its last dispatch by $T$. The calendar-time bound is
\begin{equation}
 A_\pi(T)-gT
 =\sum_{k:t_k\le T}\frac{q_k}{\Lambda}[J(q_k)-g]
     +\frac{\omega\Lambda r^2}2-gr
 \ge-\frac{g^2}{2\omega\Lambda}.
 \label{eq:analysis-calendar-bound}
\end{equation}

\begin{proof}[Proof of Theorem~\ref{thm:analysis-consistency}]
Starting immediately after a dispatch, the backlog grows from zero to $q$
in $q/\Lambda$ minutes. Its holding cost over that cycle is
$\omega q^2/(2\Lambda)$. Adding $C(q)$ and dividing by the cycle length
proves \eqref{eq:analysis-consistency}. Existence and compactness follow
from \eqref{eq:proof-growth-lower}.

Consider any admissible policy and a calendar time $T$. Admissibility
ensures that the number $N=N(T)$ of dispatches by $T$ is finite. Write
$r=T-t_N$, taking $t_0=0$ if $N=0$. Each complete cycle contributes
$(q_k/\Lambda)J(q_k)$; the unfinished cycle contributes
$\omega\Lambda r^2/2$. Subtracting $gT$ gives
\eqref{eq:analysis-calendar-bound}, because all completed-cycle excesses
are nonnegative and
\[
 \frac{\omega\Lambda r^2}2-gr
 =\frac{\omega\Lambda}2
        \left(r-\frac g{\omega\Lambda}\right)^2
       -\frac{g^2}{2\omega\Lambda}.
\]
This uniform lower bound proves the limiting inferior claim. If a threshold
$b\in\mathcal B$ is repeated, all complete-cycle excesses vanish and
$0\le r<b/\Lambda$. The residual term divided by $T$ tends to zero, so
the complete calendar-time average converges to $g$. Multiplication by
the positive constant $\Lambda$ preserves the static minimizer set.
\end{proof}

\subsection{A verified average-cost QVI on the reduced state space}
\label{subsec:qvi}

Normalize the relative value by $V(0)=0$. For $q>0$, the intervention
operator is $\mathcal MV(q)=C(q)+V(0)$. The average-cost
quasi-variational inequality is
\begin{equation}
 \min\left\{\omega q-g+\Lambda V'(q),\ C(q)-V(q)\right\}=0.
 \label{eq:analysis-qvi}
\end{equation}
This is complementarity notation: the waiting cost-rate residual and
the intervention cost residual are both nonnegative, with at least one
zero. At $q=0$
dispatch is unavailable, and the waiting equation supplies the boundary
condition $\Lambda V'(0+)=g$. The link with impulse control is the
intervention operator; QVI terminology alone does not supply verification
or an optimal policy \cite{LionsPerthame1986,Belak2017}.

In the construction below, $F$ denotes an excess-cost function, distinct
from the city routing inverse $F_k$ in Section~\ref{sec:capacity}.

\begin{theorem}[Global relative-value construction and verification]
\label{thm:analysis-qvi}
Let $g$ be as in Theorem~\ref{thm:analysis-consistency} and define
\begin{align}
 F(q)&=C(q)+\frac{\omega q^2}{2\Lambda}-\frac{gq}{\Lambda},
 &m(q)&=\inf_{z\ge q}F(z),\qquad q>0,
 \label{eq:analysis-tail-envelope}\\
 V(q)&=\frac{gq-\omega q^2/2}{\Lambda}+m(q),
 &V(0)&=0.
 \label{eq:analysis-value}
\end{align}
The tail infima are attained. The function $V$ is locally absolutely
continuous, satisfies \eqref{eq:analysis-qvi} almost everywhere on
$(0,\infty)$, and has $V'(0+)=g/\Lambda$. Along every admissible policy,
\begin{equation}
 A_\pi(T)-gT\ge V(Q(0))-V(Q(T)).
 \label{eq:analysis-verification-bound}
\end{equation}
Moreover $V(q)=O(1+q)$ in absolute value. On every sequence $T_j\uparrow\infty$
with bounded $A_\pi(T_j)/T_j$, the transversality term
$V(Q(T_j))/T_j$ tends to zero. This verifies the same average-cost lower
bound and the optimal repeated thresholds.
\end{theorem}

The tail envelope is needed for a statement on all initial backlog states.
The cycle minimizer alone does not show that immediately dispatching at
every state above its smallest value is an optimal first action. Equivalently,
the constructed relative value is the first-dispatch excess-cost envelope
\begin{equation}
 V(q)=\inf_{z\ge q}\left\{C(z)
       +\frac{\omega(z^2-q^2)}{2\Lambda}
       -\frac{g(z-q)}{\Lambda}\right\}.
 \label{eq:analysis-first-dispatch}
\end{equation}
If $b_*:=\min\mathcal B$, then $m(q)=0$ on $[0,b_*]$. At a differentiable
interior optimum $b_*$, value matching and the stationary first-order
condition give $V(b_*)=C(b_*)$ and
$C'(b_*)=(g-\omega b_*)/\Lambda$. At a load or capacity kink these
identities must be interpreted using the available one-sided derivatives;
classical smooth fit is not imposed. No uniqueness of a QVI solution or
general viscosity comparison theorem is asserted here.

\begin{proof}[Proof of Theorem~\ref{thm:analysis-qvi}]
The definition of $g$ gives
\[
 F(q)=\frac q\Lambda[J(q)-g]\ge0,
 \qquad \{q>0:F(q)=0\}=\mathcal B.
\]
The function $F$ is locally Lipschitz on $(0,\infty)$, has the finite
right limit in \eqref{eq:proof-zero-limit}, and tends to infinity as
$q\to\infty$. Its tail infima are therefore attained. It is not
necessary, and generally not correct, to assert that $F$ is Lipschitz
at zero: BHH terms may contribute $\sqrt q$ to $C(q)$.

The tail minimum $m$ is nondecreasing. It is also locally Lipschitz on
$(0,\infty)$. To see this, take a compact interval $[r,R]$ with $r>0$.
Coercivity of $F$ allows all tail minimizers for arguments in $[r,R]$
to be chosen in a fixed compact interval $[r,Z]$. Let $L_F$ be a
Lipschitz constant of $F$ there. For $r\le x<y\le R$, if a minimizer
for $m(x)$ lies at or above $y$, then $m(x)=m(y)$. Otherwise choose
one $z_x\in[x,y)$; since $m(y)\le F(y)$,
\[
 0\le m(y)-m(x)\le F(y)-F(z_x)\le L_F(y-x).
\]
At every point with $F(q)>m(q)$, a tail minimizer lies strictly to its
right. Continuity then makes $m$ constant on a neighborhood of $q$:
a small change in the lower tail boundary neither removes that minimizer
nor admits a smaller value. Hence $m'=0$ almost everywhere on the strict
noncontact region.

Let $b_* =\min\mathcal B>0$. Nonnegativity of $F$ and $F(b_*)=0$ imply
$m(q)=0$ for $0<q\le b_*$. Set $m(0)=0$. The resulting $V$ is a smooth
polynomial near zero, is locally absolutely continuous elsewhere, and
has $V(0)=0$ and $V'(0+)=g/\Lambda$. Almost everywhere on $(0,\infty)$,
\[
 C(q)-V(q)=F(q)-m(q)\ge0,
 \qquad
 \omega q-g+\Lambda V'(q)=\Lambda m'(q)\ge0.
\]
If the first quantity is positive, local constancy of $m$ makes the
second zero. If it is zero, complementarity already holds. This proves
the QVI in the stated almost-everywhere sense. Algebraically expanding
the tail-infimum definition also gives
\eqref{eq:analysis-first-dispatch}.

Integrate the waiting inequality along each uninterrupted drift segment.
At an intervention from backlog $q$ to zero, use
$C(q)\ge V(q)-V(0)$. Summing over the finitely many interventions
before $T$ telescopes all intermediate values of $V$ and proves
\eqref{eq:analysis-verification-bound}. This argument is valid for
locally absolutely continuous $V$, since backlog grows linearly on each
drift segment and thus spends zero time in the exceptional set of
nondifferentiability.

To verify transversality rather than presume it, first use immediate
dispatch in \eqref{eq:analysis-first-dispatch} to obtain $V(q)\le C(q)$.
The same formula and $C\ge0$ imply
\[
 V(q)\ge \inf_{z\ge q}
      \int_q^z\frac{\omega u-g}{\Lambda}\,du
      \ge-\frac{g^2}{2\omega\Lambda}.
\]
Together with \eqref{eq:proof-growth-upper}, this gives
$|V(q)|\le K_1+K_2q$ for finite constants. At any calendar time, the
unfinished cycle alone has incurred holding cost
$\omega Q(T)^2/(2\Lambda)$. All other costs are nonnegative. Therefore,
on any sequence with bounded $A_\pi(T_j)/T_j$,
\[
 Q(T_j)=O(\sqrt{T_j}),\qquad
 \frac{|V(Q(T_j))|}{T_j}\longrightarrow0.
\]
If the limiting inferior average cost is finite, choose such a sequence
realizing it and apply \eqref{eq:analysis-verification-bound}; if it is
infinite, the lower bound is immediate. For a repeated optimal threshold
the backlog is bounded, all relevant inequalities are equalities along
the regenerative path, and the same bound verifies its average cost $g$.
\end{proof}

The construction also describes optimal first-dispatch excess costs from
arbitrary positive backlogs: wait until a point attaining the tail minimum
in \eqref{eq:analysis-first-dispatch}, dispatch, and then repeat any
threshold in $\mathcal B$. For states above $b_*$, immediate dispatch
requires an additional shape property of $F$; for example, it follows if
$F$ is nondecreasing on $[b_*,\infty)$. Merely minimizing $J$ at $b_*$
does not establish that property. At a differentiable interior minimizer,
$F'(b_*)=0$ gives the derivative statement above. At a
nondifferentiable local minimum the valid one-sided statement is
$F'_-(b_*)\le0\le F'_+(b_*)$, when these derivatives exist.

\subsection{Unit Poisson arrivals within the integer-threshold class}
\label{subsec:analysis-limits}

There is one useful exact stochastic comparison. Let arrivals instead be a
unit-jump Poisson process of rate $\Lambda$, with its natural filtration,
and restrict decisions to dispatching all shipments immediately when the
backlog reaches a fixed positive integer $B$. Retain the same $C(B)$ and
linear holding cost. Renewal accounting gives the following result.

\begin{proposition}[A common integer-threshold comparison]
\label{prop:analysis-poisson}
For these Poisson threshold policies,
\begin{equation}
 J_P(B)=\frac{\Lambda C(B)}B+\frac{\omega(B-1)}2
       =J(B)-\frac\omega2,\qquad B\in\{1,2,\ldots\}.
 \label{eq:analysis-poisson}
\end{equation}
Consequently the Poisson and deterministic-fluid problems have the same
minimizing thresholds when both are restricted to this same integer set.
Their optimal cost rates differ by $\omega/2$.
\end{proposition}

This restricted result is not an optimality theorem over all adapted
stochastic intervention policies. The Poisson waiting generator is
$\Lambda[V(n+1)-V(n)]$, rather than $\Lambda V'(q)$; a general stochastic
verification requires integrability and expected transversality conditions.
The continuous-wave and integer-wave domains also need not have the same
optimizer. For $B>1$, the time required to collect the remaining Poisson
arrivals has no deterministic upper bound, whereas $B=1$ dispatches at
arrival and creates no consolidation wait.

\begin{proof}[Proof of Proposition~\ref{prop:analysis-poisson}]
For a fixed integer $B\ge1$, the dispatch cycle contains exactly $B$
independent exponential interarrival times of mean $1/\Lambda$. Its
mean duration is $B/\Lambda$. The backlog spends expected time
$1/\Lambda$ in each state $0,1,\ldots,B-1$, so expected holding cost
per cycle is
\[
 \frac\omega\Lambda\sum_{n=0}^{B-1}n
       =\frac{\omega B(B-1)}{2\Lambda}.
\]
The service cost per cycle is $C(B)$. The renewal-reward ratio is
therefore \eqref{eq:analysis-poisson}. The difference from the fluid
criterion is the same constant for every admissible integer threshold,
which preserves its minimizer set. Both integer minima exist, since
the holding term tends to infinity with $B$. This calculation proves
the claim within the stated threshold class and makes no comparison
with other adapted intervention rules.
\end{proof}

The restriction to a common action domain is substantive. For example,
in a separate setup-cost illustration with $C(q)=9/8$ and
$\Lambda=\omega=1$, the continuous-fluid optimum is $q=3/2$ with cost
rate $3/2$. The same fluid criterion restricted to positive integers
has its unique minimum at $B=2$, with cost rate $25/16$. The Poisson
threshold criterion also chooses $B=2$ and has cost rate $17/16$.
Thus equality of mean arrival rates alone implies neither equality
with the continuous optimizer nor equality of cost rates. This example
illustrates the action-domain issue; its constant setup cost is not
the BHH service cost in \eqref{eq:analysis-dispatch-cost}.

